\documentclass{situationist-preprint}

\title{Kripke's Contingent A Priori Reconsidered}
\msid{EPI-12-b}
\msfamily{EPI-12}
\msdiscipline{Epistemology}
\mstype{Companion manuscript}
\mskeywords{contingent a priori; reference fixing; rigid designation; Kripke; Evans; conceptual grounding}

\begin{document}
\maketitle

\begin{abstract}
Kripke argued that someone who fixes the reference of \enquote{one meter} as the length of a stick at a time knows a priori that the stick is one meter long at that time, although this is contingent. EPI-12-a holds that a proposition can be known a priori if and only if it is conceptually grounded. We argue that Kripke's verdict is correct and that his cases fall within the scope EPI-12-a allows: the stipulation introduces a concept with a constitutive rule from which the proposition is derived, and the resulting belief is safe. The knowledge so obtained is superficially contingent and deeply uninformative, so the cases do not enlarge a priori knowledge of the world beyond what EPI-12-a allows.
\end{abstract}

\section{The cases}

\citet{Kripke1980} considered a speaker who fixes the reference of \enquote{one meter} by the description \enquote{the length of stick $S$ at $t_0$}, using the description to fix the reference of a rigid designator rather than to give its meaning. In other possible worlds $S$ might have been longer or shorter at $t_0$, and \enquote{one meter} designates the same length there, so \enquote{$S$ is one meter long at $t_0$} expresses a contingent truth. Yet the speaker seems to know it without further investigation. \citet{Evans1979} gave a parallel case with a descriptive name, \enquote{Julius}, introduced as the name of whoever invented the zip.

EPI-12-a holds that there is a priori knowledge, and that a proposition can be known a priori if and only if it is conceptually grounded, that is, derivable from no premises by rules constitutive of its concepts (its Definition~2.3 and Theorem~3.3). Nothing in that thesis requires a priori truths to be necessary. The question is whether Kripke's proposition is conceptually grounded and whether the speaker's belief in it meets the safety condition.

\section{The verdict}

\begin{proposition}\label{p:ground}
For the speaker who introduces \enquote{one meter} by Kripke's stipulation, the proposition that $S$, if it has a length at $t_0$, is one meter long at $t_0$ is conceptually grounded.
\end{proposition}

\begin{proof}
The stipulation introduces a concept, \textsc{one meter}, possession of which includes a disposition to infer from \enquote{$x$ is the length of $S$ at $t_0$} to \enquote{$x$ is one meter}. That disposition is constitutive of the concept in the sense of Definition~2.2 of EPI-12-a. From no premises, the conditional that if $S$ has a length at $t_0$ then that length is one meter follows by this rule and logic.
\end{proof}

Experience figures in the case: the speaker must see or handle $S$ to refer to it. By Definition~2.1 of EPI-12-a, experience that only makes possible the possession of the concepts involved does not prevent justification from being a priori. That is exactly the role the speaker's perception of $S$ plays.

\begin{proposition}\label{p:safe}
The speaker's belief, so derived, is safe; hence, by Theorem~3.3 of EPI-12-a, it is a priori knowledge.
\end{proposition}

\begin{proof}
By Assumption~2.4 of EPI-12-a the constitutive rules of a non-defective concept are truth-preserving, the concept's semantic value being whatever makes them so; the concept introduced by the stipulation is non-defective, since its rule is satisfied by the length $S$ actually has. In close cases in which the speaker believes on the same basis, she introduces the concept by the same stipulation, and the proposition she believes is true whatever length $S$ has there. This is Lemma~3.1 of EPI-12-a applied to the case.
\end{proof}

The proposition believed is contingent, and the belief is safe, because in the close cases in which $S$ has a different length the speaker's stipulation fixes a different length as the referent, and the proposition she then believes is a different one, also true. Safety concerns the propositions believed on the same basis (Principle~P2 of EPI-12-a), not the proposition actually believed evaluated at other worlds; that is why it does not require necessity.

\section{What the cases do not show}

\citet{Evans1979} distinguished superficial from deep contingency. A sentence is superficially contingent if, embedded under a modal operator, it yields a falsehood; it is deeply contingent if what makes it true is a fact about the world that a speaker who understands it may fail to know. Kripke's sentence is superficially contingent but not deeply contingent: understanding it guarantees its truth. Propositions~\ref{p:ground} and~\ref{p:safe} explain this. The knowledge is obtained by derivation from a constitutive rule, and the rule settles nothing about the world except that, if $S$ has a length, the length is called one meter.

\begin{objection}
The speaker does not know which length one meter is, and so does not know of any length that $S$ has it; she knows only a truth about her word \citep{Donnellan1977}.
\end{objection}

\begin{reply}
She knows a proposition about a length, the one the stipulation fixes, under a concept whose only constitutive rule identifies it by $S$. She does not know how that length compares with others measured in other ways, which is knowledge of the world and is a posteriori. The objection correctly identifies the limits of her knowledge; it does not show that she lacks the knowledge Proposition~\ref{p:safe} attributes.
\end{reply}

\section{Conclusion}

Kripke's examples are genuine cases of contingent a priori knowledge. They fall under the conceptual truths of Proposition~3.6 of EPI-12-a, being derived from constitutive rules of concepts introduced by stipulation, and their contingency is superficial. They confirm that the a priori is fixed by conceptual grounding rather than by necessity, and they add nothing to a priori knowledge of how the world is.

\begin{thebibliography}{9}\small
\bibitem[Donnellan(1977)]{Donnellan1977} Donnellan, K.~S. (1977). The contingent a priori and rigid designators. \emph{Midwest Studies in Philosophy}, 2, 12--27.
\bibitem[Evans(1979)]{Evans1979} Evans, G. (1979). Reference and contingency. \emph{The Monist}, 62(2), 161--189.
\bibitem[Kripke(1980)]{Kripke1980} Kripke, S.~A. (1980). \emph{Naming and Necessity}. Harvard University Press.
\end{thebibliography}

\end{document}
